\section{Fine-tuning specializes a guard: a large represented gain that does not transfer (Act~I)}
\label{sec:actI}

The first act answers the most basic question a practitioner has after fine-tuning a guard:
\emph{what did the fine-tune actually buy, and does it hold up off the training distribution?} The
leaderboard answer --- a single higher number --- silently blends two very different things: doing
better on benchmark \emph{sources the guard trained on} (``represented''), and doing better on
datasets it \emph{never saw} (``transfer''). Act I separates them with a strictly paired,
same-checkpoint measurement and finds a clean, repeatable pattern: supervised fine-tuning (SFT)
delivers a large, uniform represented-source gain and essentially no average transfer gain --- the
guard \emph{specializes}.

\begin{background}{What ``paired'' buys us, and why it matters}
Most guard papers compare \emph{different} models: model $X$'s guard against model $Y$'s guard. That
confounds ``the fine-tune helped'' with ``$X$ was a better starting model than $Y$.'' We instead hold
the checkpoint fixed and compare each guard \emph{against its own base}, on the identical rows, with
the identical single-token score head ($s(x)=z_{\text{unsafe}}-z_{\text{safe}}$; see
\Cref{sec:background}) and the identical tie-aware macro-AP. The reported quantity is therefore a
\emph{within-checkpoint change}, $\Delta = \text{SFT} - \text{base}$, not a cross-model gap. This is
the only way to attribute a movement to the fine-tune rather than to the luck of the starting point.
\end{background}

\subsection{The paired base-to-SFT design}
\label{sec:actI-design}

\paragraph{Panel and manifest.} All four checkpoints (Qwen2.5-1.5B, SmolLM2-1.7B, SmolLM3-3B,
Qwen3-4B) are fine-tuned from the identical frozen manifest of $1{,}200$ rows: $400$ from each of three
represented sources --- \code{toxicchat}, \code{prompt\_injections}, and
\code{jailbreak\_classification} --- balanced $200$ safe / $200$ unsafe within each source. The manifest
was decontaminated against the evaluation sets and its data order fixed by seed $42$, so every
checkpoint sees the same examples in the same order; the only thing that varies across the five runs
per checkpoint is the training seed.

\paragraph{Recipe.} Each guard is a LoRA adapter (\code{r=32}, \code{alpha=64}, \code{dropout=0.05})
attached to the attention and MLP projections (\code{q,k,v,o,gate,up,down}), trained for $300$ steps at
learning rate \code{2e-4} on a cosine schedule with $3\%$ warmup, effective batch size $4$, and maximum
sequence length $1024$. We run five seeds ($42$--$46$) per checkpoint, giving $20$ SFT guards in total.
Only the adapter is trained; the base weights are frozen, so ``base'' means the same model \emph{before}
the guard adapter, evaluated with the identical scoring head.

\paragraph{Two evaluation regimes.} Every guard and its base are scored in two disjoint regimes.
\emph{Represented} = held-back rows from the three training sources (same sources, unseen rows);
\emph{transfer} = four datasets whose rows were never used in training --- \code{jailbreakbench},
\code{xstest}, \code{wildguardtest}, and \code{wildjailbreak}. We additionally probe two single-class
stress sets that isolate the two failure directions: \code{OR-Bench} (benign prompts, to measure
over-refusal) and \code{HarmBench} (hard attack prompts, to measure missed catches).

\paragraph{Estimand and intervals.} We summarise each regime by macro-AP (average precision averaged
with equal weight across benchmarks, so a large benchmark cannot dominate) and put a two-sided $95\%$
interval on each change with a \emph{paired hierarchical bootstrap} that keeps every adapter tied to its
own base and respects the checkpoint$\to$seed grouping. The estimand is deliberately narrow: it is the
change \emph{conditional on this fixed, purposively-chosen four-checkpoint panel and this manifest},
which was inspected during development. These are estimation intervals, not significance tests, and the
result is \emph{retrospective, not confirmatory} --- no causal or population claim is licensed.

\subsection{The primary result: a uniform represented gain, a split transfer effect}
\label{sec:actI-primary}

\Cref{tab:primary} reports, per checkpoint, the untuned-base and mean-SFT macro-AP in each regime, the
paired base-to-SFT change with its interval, and the fixed-panel aggregate.

\begin{table}[htbp]\centering\small
\caption{Act I --- fixed-panel result: base and mean-SFT macro-AP by regime, paired base-to-SFT deltas
with two-sided $95\%$ hierarchical-bootstrap intervals, and the fixed-panel aggregate. Descriptive under
the locked precision-focused mode; conditional on this panel.}
\label{tab:primary}
\resizebox{\textwidth}{!}{\input{generated/tab_primary_gen}}
\end{table}

\paragraph{Represented: everyone reaches the same ceiling.} The represented-source gain is large and,
strikingly, \emph{uniform in its destination}: regardless of where the base started, every SFT guard
lands at macro-AP $\approx 0.98$. SmolLM2-1.7B climbs from a weak $0.4524$ to $0.9806$; Qwen2.5-1.5B
from $0.6334$ to $0.9878$; SmolLM3-3B from $0.6621$ to $0.9751$; and Qwen3-4B, already strong at
$0.8855$, to $0.9837$. Because the finish line is shared, the \emph{size} of the represented gain is
essentially dictated by how far below the ceiling the base sat --- a $+0.5282$ jump for SmolLM2 versus
only $+0.0981$ for Qwen3-4B. The fixed-panel aggregate is $\RepDelta$
$[\RepCILower, \RepCIUpper]$: SFT reliably teaches each model to rank the training sources' held-back
rows near-perfectly. Mechanistically this is unsurprising --- the manifest \emph{is} those sources, so
the adapter learns their label conventions and surface cues.

\paragraph{Transfer: a small average that hides opposite signs.} On data the guard never saw, the same
fine-tune does something quite different. The aggregate transfer change is only $\TransferDelta$
$[\TransferCILower, \TransferCIUpper]$ --- close to zero and slightly negative. But this panel average
is a mirage: it pools \emph{opposing} per-checkpoint effects. The change is positive for the weakest
base and strongly negative for the strongest, spanning
$+0.0400$ $[+0.0003, +0.0776]$ for SmolLM2-1.7B, $-0.0389$ $[-0.0829, +0.0062]$ for Qwen2.5-1.5B,
$-0.0869$ $[-0.1114, -0.0613]$ for SmolLM3-3B, and $-0.1499$ $[-0.1963, -0.1050]$ for Qwen3-4B. The
ordering tracks base transfer strength (\Cref{tab:primary}): the models with the most transfer skill to
begin with (SmolLM3 at $0.9102$, Qwen3-4B at $0.9438$) give the most of it back, while the model that
had the least ($0.7904$) is the only one that improves. In other words, SFT does not add a portable
notion of ``unsafe''; it reshapes the score toward the training sources, and for a model that already
generalised well that reshaping is a net loss off-distribution.

\begin{takeaway}{Why the $\TransferDelta$ transfer average is misleading}
It averages effects that point in opposite directions --- SmolLM2 $+0.0400$ against Qwen3-4B $-0.1499$.
The panel-level number is small; the panel itself is \emph{not} uniform, and reporting only the
aggregate would erase the most important finding (who loses, and how much). Read the per-checkpoint
column, not the bottom row.
\end{takeaway}

\subsection{Where the transfer losses live: a per-benchmark decomposition}
\label{sec:actI-perbench}

Aggregating over benchmarks can also hide \emph{which} data the guard forgets how to rank. The upper
block of \Cref{tab:sensitivity} decomposes the change benchmark by benchmark. On the represented side
every source rises sharply --- \code{toxicchat} $+0.1880$, \code{prompt\_injections} $+0.3704$, and
\code{jailbreak\_classification} $+0.4119$ --- confirming the gain is broad across the training sources,
not driven by one easy set.

On the transfer side the losses are \emph{concentrated on the jailbreak-style adversarial sets}:
\code{jailbreakbench} $-0.0776$, \code{wildjailbreak} $-0.0792$, and \code{wildguardtest} $-0.0669$ all
fall by a comparable amount, while \code{xstest} --- which contrasts genuinely unsafe prompts against
benign look-alikes rather than adversarial attacks --- barely moves at $-0.0120$. The pattern is
coherent with the mechanism above: the training sources over-represent a particular flavour of unsafe
text, so after tuning the guard ranks \emph{those} confidently but loses discrimination precisely on the
adversarial, distribution-shifted prompts that a deployed guard most needs to catch. Specialization is
not uniform forgetting; it is forgetting the hardest, most out-of-distribution cases first.

\subsection{The specialization plane: 15 of 20 guards specialize}
\label{sec:actI-plane}

The clearest single view of Act I is the \emph{specialization plane} (\Cref{fig:plane}), which plots
each of the $20$ (checkpoint, seed) guards by its represented change (horizontal) against its transfer
change (vertical). Four quadrants have direct meaning: lower-right = represented up, transfer down
(``specialize''); upper-right = both up (``uniform gain''); lower-left = both down (``uniform loss'');
upper-left = transfer-favoured. \SpecializationSeedCount{} of \TotalSeedCount{} guards land in the
specialize quadrant, \UniformGainSeedCount{} in uniform gain, and \emph{none} in uniform loss or the
transfer-favoured quadrant. So specialization is the dominant --- but not universal --- outcome: the
five uniform-gain points are four of the five SmolLM2 seeds \emph{plus} Qwen2.5's seed~42 --- concentrated
on the two checkpoints with the weakest base transfer, which have the most room to improve there (SmolLM2's
remaining seed, 46, narrowly specializes at $-0.001$). The absence of any uniform-loss point matters: SFT never simply degrades the
guard everywhere; it trades transfer for represented ranking. Per-seed values behind the plane are
tabulated in \Cref{tab:seed-values}, and show the effect is stable within each checkpoint (e.g. every
Qwen3-4B seed is a substantial transfer loss, every SmolLM2 seed a represented gain).

\begin{figure}[htbp]\centering
\includegraphics[width=0.6\textwidth]{specialization_plane.pdf}
\caption{Each point is one (checkpoint, seed): horizontal axis is the represented-source macro-AP
change, vertical axis the transfer change. \SpecializationSeedCount{}/\TotalSeedCount{} land in the
lower-right specialization quadrant --- the visual signature of Act I.}
\label{fig:plane}
\end{figure}

\begin{table}[htbp]\centering\scriptsize\setlength{\tabcolsep}{2.5pt}
\caption{Observed represented-source and transfer macro-AP change for every SFT seed (Act~I),
generated from the same validated score matrix as \Cref{tab:primary}. Every Qwen3-4B seed is a substantial
transfer loss and every SmolLM2-1.7B seed a represented gain --- the per-seed stability behind the
specialization plane (\Cref{fig:plane}).}
\label{tab:seed-values}
\input{generated/tab_seed_values_gen}
\end{table}

\subsection{At a deployable operating point}
\label{sec:actI-op}

Ranking is threshold-free, but deployment is not: at some point you must pick a cutoff and turn scores
into \code{block}/\code{allow} decisions. To see the deployment consequence of specialization we select
each guard's threshold to hit a $5\%$ false-positive-rate (FPR) target on a calibration split, then read
off the realized rates. The lower blocks of \Cref{tab:sensitivity} report them.

\paragraph{Represented recall soars.} On the represented sources, catching the unsafe prompts (TPR)
jumps from \RepBaseTPRPct\% (base) to \RepSFTTPRPct\% (SFT) --- the operating-point face of the same
$\approx 0.98$ ranking gain --- while the represented false-alarm rate stays low and even edges down
(\RepBaseFPRPct\% to \RepSFTFPRPct\%). If you only measured on the training sources, the guard would
look strictly and dramatically better.

\paragraph{Transfer pays for it.} Off-source the picture inverts. The transfer benchmark-macro FPR
climbs from \TransferBaseFPRPct\% (base) to \TransferSFTFPRPct\% (SFT) --- the tuned guard raises nearly
twice as many false alarms on data it never trained on --- and the pooled-negative FPR blows past target
from \TransferBasePooledFPRPct\% to \TransferSFTPooledFPRPct\%, so the calibrated threshold simply does
not transfer. Transfer recall rises only modestly (\TransferBaseTPRPct\% to \TransferSFTTPRPct\%),
nowhere near the represented jump. The single-class stress sets make the cost concrete: on the hard
\code{HarmBench} attacks the tuned guard's recall \emph{falls} from \HarmBenchBaseRecallPct\% to
\HarmBenchSFTRecallPct\% --- it catches less of the very content a guard exists to stop --- while
benign over-refusal on \code{OR-Bench} is essentially flat (\ORBenchBaseFPRPct\% to \ORBenchSFTFPRPct\%),
so the extra false alarms are an off-distribution phenomenon, not a blanket increase in caution.

\begin{table}[htbp]\centering\small\setlength{\tabcolsep}{6pt}
\caption{Act I --- per-benchmark paired deltas (upper) and the calibration-targeted $5\%$ FPR operating
point plus single-class stress diagnostics (lower). Generated from the lock-bound scores.}
\label{tab:sensitivity}
\input{generated/tab_sensitivity_gen}
\end{table}

\begin{takeaway}{The deployment cost behind the ranking story}
Read together, the operating point says the specialized guard \emph{catches more of what it trained on,
raises more false alarms on what it didn't, and misses more hard attacks}. A leaderboard number computed
on represented sources would advertise only the first of these three. The threshold that looked safe in
calibration is not the threshold you get in the field --- a theme Acts~II and~III return to.
\end{takeaway}

\subsection{The attractor: post-SFT scores are benchmark-fixed, so ``stronger bases specialize more'' is arithmetic}
\label{sec:actI-coupling}

The cleanest way to see what Act~I does and does not show is to notice a regularity that runs underneath
\Cref{tab:primary}: \textbf{fine-tuning behaves like an attractor}. No matter where a checkpoint starts,
SFT pulls its ranking to nearly the \emph{same} score on a given benchmark. Represented-source AP after SFT
is $0.982 \pm 0.005$ across the four checkpoints; transfer AP after SFT is $0.807 \pm 0.024$ --- both far
tighter than the \emph{base} spread they came from ($0.178$ and $0.073$ respectively; the transfer variance
shrinks about \textbf{9-fold}). \Cref{fig:attractor} (left) shows the collapse: four widely spread base
scores, one narrow post-SFT band.

\begin{figure}[htbp]\centering
\includegraphics[width=0.99\linewidth]{fig_attractor.pdf}
\caption{\textbf{The fine-tuning attractor.} \emph{Left:} on the transfer benchmarks the four base
checkpoints span $0.79$--$0.94$ macro-AP, but after SFT they collapse into a narrow band around a
benchmark-fixed endpoint $A_T\approx0.81$. \emph{Right:} because every checkpoint lands near that same
endpoint, the paired change $\Delta=\text{SFT}-\text{base}$ is forced onto a line of slope $-1$ (green);
the four points sit on it (fitted slope $-1.05$, $R^2=0.91$). ``Stronger bases specialize more'' is
therefore arithmetic, not a discovered behavioral law.}
\label{fig:attractor}
\end{figure}

\paragraph{Why this reframes ``stronger bases specialize more.''} Once the endpoint is (approximately)
fixed, the change we plot is pure subtraction: $\Delta = \text{endpoint} - \text{base}$. A stronger base
then \emph{must} show a smaller $\Delta$ --- not because strong models are intrinsically fragile, but
because they start closer to the ceiling. This is testable. Regressing $\Delta$ on the
base score across the four checkpoints gives slope $-0.99$ (represented, $R^2=0.999$) and $-1.05$ (transfer,
$R^2=0.91$) --- the $-1$ the attractor \emph{predicts before any fitting}. One caveat on reading those
numbers: the represented $R^2=0.999$ is close to \emph{definitional}. When the endpoint band is as tight as
it is on represented sources ($\pm0.005$), $\Delta=\text{endpoint}-\text{base}$ is arithmetically almost a
slope-$-1$ line in the base score, so a near-perfect fit is nearly forced and is not independent
confirmation of anything. The transfer fit, where the endpoint band is several times wider, is the
load-bearing test --- and it still lands at slope $-1.05$. Either way the eye-catching ordering
(SmolLM2 $+0.040$ down to Qwen3-4B $-0.150$) carries no behavioral content beyond the endpoints themselves.

\paragraph{Where the real content lives.} Two places, both benchmark-owned rather than model-owned. First,
\emph{where the endpoints sit}: $A_T\approx0.81$ is largely a property of the \emph{(training manifest,
benchmark)} pair --- to the precision of the small residuals below, the base model has largely dropped out
of the equation, and the post-SFT score is \emph{largely determined by} the manifest and the benchmark
rather than by which checkpoint you started from. This is the strongest form of the paper's thesis:
after SFT, the benchmark and the manifest largely choose the score. Second, the \emph{small residuals} around the
endpoint (at most $0.027$ on transfer): these measure how much of a base's identity survives the fine-tune,
and are arguably a more honest per-checkpoint quantity than $\Delta$ itself. LoRA only adds a low-rank
correction in the base's own features, so full base-independence is approximate by construction --- the
residual is exactly that approximation, made visible.

\begin{background}{Why the represented ceiling is $\approx0.98$, not $1.0$ --- and why that helps the thesis}
A perfect ranker cannot score AP $=1$ on a benchmark whose \emph{labels} are slightly noisy: a small rate
$\eta$ of ambiguous or mislabeled rows caps the achievable AP at roughly $1-1.3\,\eta$ (the coefficient
$1.3$ is an order-of-magnitude rule of thumb, not derived here --- the exact factor depends on prevalence
and on where mislabels fall in the ranking). The observed common
ceiling $0.982$ is what you would see at $\eta\approx1.5\%$ label noise --- i.e.\ the ceiling would then be a
property of the benchmark's \emph{labels}, not of the models. This is a \emph{hypothesis}, not a measured
result: it is checkable by auditing the top-ranked ``false positives'' of the SFT guards --- under this
reading most should be mislabels --- but we have not run that audit, so the label-noise explanation of the
ceiling remains conjectural. Either way, the empirical fact we rely on is only that a common ceiling exists
and is shared across checkpoints; the benchmark, not the model, sets the top of the scale.
\end{background}

\begin{takeaway}{How much to trust the ``base strength'' ordering}
The apparent trend rests on \emph{four} checkpoints --- and because seeds within a checkpoint are highly
correlated (every Qwen3-4B seed is negative, every SmolLM2 seed but one positive), the effective sample is
$n=4$, not $20$. But the attractor turns the checkpoint-level ``specialize'' pattern into a \emph{prediction} from a
single constant: a checkpoint specializes exactly when its base transfer score exceeds the endpoint
($\text{base}_T>A_T\approx0.81$), which selects Qwen2.5, SmolLM3 and Qwen3-4B as specializers and leaves
SmolLM2 in uniform-gain --- the checkpoint-mean split we observe, and a $15/5$ seed count. The match is at
the checkpoint mean, not seed-perfect: two of the twenty seeds cross their checkpoint's boundary (Qwen2.5's
seed~42 gains $+0.008$; SmolLM2's seed~46 loses $-0.001$), so the constant predicts the aggregate split
rather than every individual seed. The sharper, falsifiable hypothesis to carry forward is
therefore \emph{endpoint invariance} itself: any new $1.5$--$4$B instruct checkpoint tuned with this recipe
should land near represented $0.98$ and transfer $0.81$ regardless of where its base started. One honest
limit on this extrapolation: our four checkpoints are only \emph{two} lineages (two Qwen, two SmolLM), so
what we have actually shown is endpoint invariance \emph{within} two families --- a prospective test
(roadmap item~7) should add unrelated lineages before the invariance is treated as recipe-general. What
Act~I establishes now is the qualitative claim, robust across all four checkpoints and $20$ seeds --- SFT
buys represented-source ranking and not, on average, transfer.
\end{takeaway}

\subsection{A recipe control: does anti-forgetting (KL-regularized) SFT preserve transfer?}
\label{sec:actI-klsft}

The specialization above was measured under \emph{one} fine-tuning recipe --- completion-only LoRA-SFT with
no explicit anchor to the base checkpoint. A fair objection is that the transfer loss might be a property of
this \emph{unregularized} recipe rather than of fine-tuning as such. The standard remedy is an
\emph{anti-forgetting} penalty that keeps the fine-tune close to the base in output space, so we add exactly
that --- a KL term on the SFT loss:
\begin{equation}
\mathcal{L} \;=\; \underbrace{\mathrm{CE}(\text{verdict})}_{\text{Act~I recipe}} \;+\;
\beta\,\mathrm{KL}\!\big(\pi_\theta(\cdot\mid x)\,\big\|\,\pi_{\text{base}}(\cdot\mid x)\big),
\label{eq:klsft}
\end{equation}
evaluated on the completion tokens, with the frozen base $\pi_{\text{base}}$ recovered from the adapter's own
disabled path (no second model in memory, no base retraining). Everything else is held at the Act~I settings
--- same manifest, seeds, panel, scorer, and represented/transfer splits --- and $\beta=0$ reproduces vanilla
SFT \emph{exactly} (the cross-entropy is unchanged), so this is a strict one-knob generalization of the Act~I
recipe rather than a different method. We sweep $\beta\in\{\KLBetas\}$ over all four checkpoints and
\KLNSeeds\ seeds, and score every adapter through the identical single-token margin used everywhere else.

% GENERATED by experiments/analyze_klsft.py -- do not edit by hand.
% PLACEHOLDER pending the KL-SFT GCP run; row values fill in once the klsft parquets are analyzed.
\begin{table}[H]\centering\footnotesize
\caption{\textbf{Anti-forgetting control (KL-regularized SFT).} Transfer and represented macro-AP for the
base, vanilla SFT ($\beta{=}0$, trained in the same environment as the KL runs), and KL-regularized SFT
(\Cref{eq:klsft}) at $\beta{=}0.5$ and $\beta{=}1.0$, over 5 seeds, scored identically to Act~I. The
transfer block answers whether a base-anchored penalty preserves the transfer that vanilla SFT gives up;
the represented block is the cost. \emph{(Values pending the running sweep.)}}
\label{tab:klsft}
\begin{tabular}{l cccc c ccc}\toprule
 & \multicolumn{4}{c}{transfer macro-AP} & & \multicolumn{3}{c}{represented macro-AP} \\
\cmidrule(lr){2-5}\cmidrule(lr){7-9}
Checkpoint & base & SFT & KL$_{.5}$ & KL$_{1}$ & & SFT & KL$_{.5}$ & KL$_{1}$ \\
\midrule
Qwen2.5-1.5B & -- & -- & -- & -- & & -- & -- & -- \\
SmolLM2-1.7B & -- & -- & -- & -- & & -- & -- & -- \\
SmolLM3-3B   & -- & -- & -- & -- & & -- & -- & -- \\
Qwen3-4B     & -- & -- & -- & -- & & -- & -- & -- \\
\bottomrule
\end{tabular}
\end{table}


\Cref{tab:klsft} reports transfer and represented macro-AP for the base, vanilla SFT ($\beta=0$, trained in
the \emph{same} environment as the KL runs so the contrast carries no hardware confound), and
KL-regularized SFT at each $\beta$. The anti-forgetting question is read from the transfer block: relative to
vanilla SFT, KL at $\beta=0.5$ changes transfer macro-AP by \KLBetaHalfTransferDelta\ on average across the
four checkpoints (at a represented cost of \KLBetaHalfRepDelta), and at $\beta=1.0$ by
\KLBetaOneTransferDelta\ (represented cost \KLBetaOneRepDelta). \KLVerdictSentence

\begin{takeaway}{What the recipe control shows}
\KLTakeaway
\end{takeaway}

\subsection{The deployment base rate also chooses the winner}
\label{sec:actI-prevalence}

One more axis quietly co-produces the score, and it is not the model at all: the \emph{prevalence} of
unsafe prompts. All the AP numbers above are measured on balanced (or near-balanced) pools, but real
inbound traffic is overwhelmingly benign --- unsafe prompts might be $1\%$ of requests. Average precision
depends on that base rate, and the dependence is exact: given a guard's ranking (its ROC), AP at any
prevalence $\pi_+$ is a fixed recomputation,
\begin{equation}
\mathrm{AP}(\pi_+) \;=\; \int_0^1 \frac{\pi_+\,u}{\pi_+\,u + (1-\pi_+)\,\mathrm{FPR}(u)}\,\mathrm{d}u,
\label{eq:ap-prevalence}
\end{equation}
where $u$ is recall and $\mathrm{FPR}(u)$ is the guard's false-positive rate at that recall --- so no new
model runs are needed, only the committed per-row scores.

\begin{figure}[htbp]\centering
\includegraphics[width=0.82\linewidth]{fig_prevalence.pdf}
\caption{\textbf{The prevalence also chooses the winner} (measured: \Cref{eq:ap-prevalence} applied to the
four base guards' committed transfer ROC, macro-averaged over the transfer sources). As unsafe prompts get
rarer, every guard's AP falls --- the strongest ranker (Qwen3-4B) drops from $0.94$ at balance to
$\approx0.56$ at $1\%$ prevalence, the weakest (Qwen2.5-1.5B) from $0.82$ to $0.11$ --- \emph{and} the
lower half re-orders (see text).}
\label{fig:prevalence}
\end{figure}

Two consequences follow, both recomputed exactly from the committed scores. First, the balanced AP is an
\emph{optimistic} reading of deployed precision: SmolLM3-3B's base transfer ranking scores $\mathrm{AP}=0.91$
on a balanced pool but only $\approx0.51$ at $1\%$ prevalence (\Cref{fig:prevalence}). Second --- and this is
the thesis again, now at the metric's own prior --- \textbf{low prevalence collapses and re-spaces the
ranking}. The extremes are stable (Qwen3-4B stays first, SmolLM2 and Qwen2.5 stay in the lower half at every
prevalence), but the \emph{lower two re-order}: Qwen2.5-1.5B edges SmolLM2-1.7B at balance ($0.82$ vs.\ $0.79$)
yet falls behind it once positives are rare ($0.11$ vs.\ $0.21$ at $1\%$), because a scarce positive class
re-expands exactly the low-recall precision differences the balanced pool compresses. This is a re-spacing and
a partial re-order, not a wholesale winner-flip --- but it is enough that a single balanced AP can invert two
guards' apparent order at deployment prevalence. Reporting each headline guard as a curve $\mathrm{AP}(\pi_+)$
rather than a single balanced number is a zero-cost recompute (roadmap), and it is the honest way to state a
deployment precision.

\edb{Represented AP $\RepDelta$ (LCB $\RepLCB$) vs.\ transfer $\TransferDelta$ (UCB $\TransferUCB$);
specialization in $\SpecializationSeedCount/\TotalSeedCount$ seeds (\Cref{tab:primary}). A base-anchored
KL penalty ($\beta{=}0.5$) buys back transfer (\KLBetaHalfTransferDelta{} vs.\ SFT) at a represented cost
\KLBetaHalfRepDelta{} (\Cref{sec:actI-klsft}).}
{Never replace a base guard on represented AP alone; compare each tune to its \emph{own} base on
represented \emph{and} held-out sets, and if you must SFT while caring about OOD, add the KL penalty.}
{Retrospective on a fixed four-checkpoint panel with an inspected manifest (overlap audit pending); a
description of this recipe on these runs, not a population or causal law.}
